\subsection{Linear problems}
\label{app:catalogue-linear}

\begin{table}[htbp]
    \centering
    \caption[The linear benchmark problems R1--R5]{\textbf{The linear benchmark problems R1--R5, as run.} Equation, constraint rows, datum point $x_s$ and reference solution. $x_z$, the root of the reference (or of its derivative) nearest the value Wu et al.\ print, set up from the reference; $P_l$, $P_l^m$, Legendre and associated Legendre functions; $v=f/\sqrt{1-x^2}$, the variable R3c and R3d are solved for. Every source enters by its Chebyshev coefficients; the heat equation's initial function is imposed at $t=-1$. Wu et al.'s printed rows, run only in the reproduction: \Cref{app:reproduction}.}
    \label{tab:catalogue-linear}
    \footnotesize
    \begin{tabular}{@{}l p{0.26\linewidth} p{0.29\linewidth} l p{0.24\linewidth}@{}}
        \toprule
        & Equation & Constraint rows & $x_s$ & Reference \\
        \midrule
        R1  & $f' + 2f = 0$ & -- (regular datum $f(-1)=1$) & $-1$ & $e^{-2(x+1)}$ \\
        R2a & $f'' + 4f' + 4f = 0$ & $f(-1) = 0$ & $0$ & $\tfrac12(1+x)e^{-2x}$ \\
        R2b & $f'' - 2f' - 3f = 0$ & $f(0) = 0$ & $\tfrac12$ & $\tfrac14\bigl(e^{3x} - e^{-x}\bigr)$ \\
        R2c & $f'' + 5f' + 400f = 0$ & $f'(0) = 0$ & $0$ & $e^{-5x/2}\bigl(\cos\omega x + \tfrac{\sqrt7}{21}\sin\omega x\bigr)$, $\omega = \tfrac{15\sqrt7}{2}$ \\
        R3a & $(1-x^2)f'' - 2xf' + 20f = 0$ & $f'(0) = 0$ & $1$ & $P_4$ \\
        R3b & $(1-x^2)f'' - 2xf' + 30f = 0$ & $f(0) = 0$ & $1$ & $P_5$ \\
        R3c & $(1-x^2)v'' - 4xv' + 28v = 0$ & $v'(0) = 0$ & $\tfrac12$ & $f = P_5^1 = -\sqrt{1-x^2}\,P_5'$ \\
        R3d & $(1-x^2)v'' - 4xv' + 40v = 0$ & $v(0) = 0$ & $\tfrac12$ & $f = P_6^1 = -\sqrt{1-x^2}\,P_6'$ \\
        R4a & $(x-1)f'' - xf' + f = (x-1)^2$ & $f'(x_z) = 0$, $x_z = -0.19598$ & $0$ & $\tfrac32e^{x} - \tfrac18x(8x+13) - 1$ \\
        R4b & $f'' + 4f' + 4f = e^{-2x}$ & $f(x_z) = 0$, $x_z = 1-\sqrt3$ & $0$ & $\tfrac12e^{-2x}(x^2 - 2x - 2)$ \\
        R4c & $f'' - 5f' + 6f = xe^{2x}$ & $f(x_z) = 0$, $x_z = 0.90542$ & $0$ & $\tfrac16\bigl(4e^{3x} - e^{2x}(3x^2 + 6x + 2)\bigr)$ \\
        R4d & $f'' + 4f' + 20f = xe^{-2x}\sin 4x$ & $f'(x_z) = 0$, $x_z = -0.74681$ & $0$ & $\tfrac1{64}e^{-2x}\bigl((x-48)\sin4x - 4(x^2+16)\cos4x\bigr)$ \\
        R5a & Laplace, $f_{xx} + f_{yy} = 0$ & $f(\pm1, y) = f(x, -1) = 0$; datum slice $f(x,1) = \sin\tfrac{\pi(x+1)}{2}$ & $(0, 1)$ & $\cos\tfrac{\pi x}{2}\,\sinh\tfrac{\pi(y+1)}{2}/\sinh\pi$ \\
        R5b & heat, $f_t = \tfrac{1}{25} f_{xx}$ & $f(t, \pm1) = 0$; datum slice $f(-1,x) = e^{4\pi^2/25}\sin2\pi x$ & $(-1, \tfrac14)$ & $e^{-4\pi^2 t/25}\sin 2\pi x$ \\
        R5c & wave, $f_{tt} = 4 f_{xx}$ & $f(t, \pm1) = 0$; datum slice $f(0,x) = \sin2\pi x$; $\partial_tf(0,x) = 0$ & $(0, \tfrac14)$ & $\cos 4\pi t\,\sin 2\pi x$ \\
        \bottomrule
    \end{tabular}
\end{table}

\subsection{Nonlinear and coupled problems}
\label{app:catalogue-nonlinear}

\begin{table}[htbp]
    \centering
    \caption[The nonlinear and coupled benchmark problems R6--R10]{\textbf{The nonlinear and coupled benchmark problems R6--R10.} Equation, data and reference. $x_s$, datum point; $x_z$, exact root nearest the printed value; $\rho$, Carleman ratio; $a$, initial amplitude; $\nu$, viscosity; $k\le M$, wavenumbers of the $M$ sine modes; $\omega_0$, initial vorticity; $J$, the advective nonlinearity; $w_0$, initial mode vector; \gls{ivp}, initial-value problem; \gls{dns}, direct numerical simulation; (truncation), the reference for the Galerkin truncation's error.}
    \label{tab:catalogue-nonlinear}
    \small
    \begin{tabular}{@{}l p{0.34\linewidth} p{0.24\linewidth} p{0.28\linewidth}@{}}
        \toprule
        & Equation & Data & Reference \\
        \midrule
        R6a & $4f'' + 2f'^2 + f = 0$ (Wu et al.\ Eq.~29) & $f'(0) = 0$, $x_s = 0$ & $1 - x^2/8$ \\
        R6b & $f'' - 2f^2 + x = 0$ (Wu et al.\ Eq.~34) & $f(x_z) = 0$, $x_z = 0.026147$ (printed $0.02615$) & boundary-value solve to $10^{-12}$ \\
        R7  & $u' = v$, $v' = -\omega^2 u$, $\omega = \pi$ & $u(-1) = 0$, $v(-1) = \omega$ & $u = \sin\omega(x+1)$ \\
        R8  & $y' = -y + y^2$, $t \in [0, 3]$ & $y(0) = \tfrac12$ ($\rho = 0.5$) & $1/(1 + e^{t})$ \\
        R9  & viscous Burgers, periodic, $u_0 = a\sin x$ & sine modes $k \le M$, $M = 4, 8, 16$; $a/\nu = 0.03$--$10$ & Galerkin \gls{ivp}; Cole--Hopf (truncation); exact lift \\
        R10a & 2-D vorticity Navier--Stokes, $\omega_0 = 2\sin x\sin y$ & $J \equiv 0$ (Taylor--Green) & exact viscous decay \\
        R10b--d & 2-D vorticity Navier--Stokes, counter-rotating vortex pair & $\|w_0\| = a$, $a/\nu = 3, 10, 30$; active triads & Galerkin \gls{ivp}; \gls{dns} (truncation); exact lift \\
        \bottomrule
    \end{tabular}
\end{table}

\paragraph{The fluid runs.}
\label{app:fluid-data}
Burgers' equation (\Cref{eq:burgers}) is run from $u_0=a\sin x$ over $T=1$ at five ratios $a/\nu=0.03$, $0.3$, $1$, $3$ and $10$ (R9a to R9e), each truncated to $M=4$, $8$ and $16$ sine modes; its Galerkin system, and the Cole--Hopf solution against which the truncation's own error is measured, are derived with the other Fourier--Galerkin operators in \Cref{app:fluids}.

\paragraph{NS Initial conditions.} The Taylor--Green vortex $\omega_0=2\sin x\sin y=\cos(x-y)-\cos(x+y)$ occupies two modes with $|\mathbf k|^2=2$, whose coupling vanishes, so $\omega=\omega_0e^{-2\nu t}$ exactly in any truncation that holds them, and in the equation itself. Its data lie in a subspace that $F_1$ preserves and on which $F_2$ vanishes, so its lift is exact at every order (\Cref{prop:exact-lift}): R10a validates the lift, the time axis and the preparation at a known answer. The counter-rotating pair is $\pm e^{-|\mathbf x-\mathbf c|^2/(2\sigma^2)}$ with $\sigma^2=\tfrac14$ at $\mathbf c=(3\pi/4,\pi)$ and $(5\pi/4,\pi)$, periodised and projected exactly onto the modes: it has no cosine part, and its sine coefficients are proportional to $(-1)^{k_x+k_y+1}\sin(\pi k_x/4)\,e^{-|\mathbf k|^2/8}$, so modes with $k_x\equiv0\pmod4$ have zero datum. It spans several $|\mathbf k|$, so its triads are active. It is scaled to $\|w_0\|=a$ and run at $a/\nu=3$, $10$ and $30$ (R10b to R10d) over $T=1$, and at $a/\nu=30$ also over $T=1/8$. The reference for the lift and the time axis is the Galerkin system's \gls{ivp}, with its exact sparse Jacobian; the truncation's own error is measured against a pseudo-spectral \gls{dns} of \Cref{eq:navier-stokes} from the same data, whose slope inside the truncation is the Galerkin system's.

\subsection{Run Options}
\label{app:run-options}

A problem's data enter as the four kinds of row of \Cref{def:standard-form}, listed with the problems that use them in \Cref{tab:constraint-kinds}.

\begin{table}[htbp]
    \centering
    \caption[The four kinds of constraint row]{\textbf{Every constraint row either leaves the solution's scale free or ties data to it.} The four rows that carry a problem's data, grouped by that role. Each enters the stacked residual at weight one, is the same in both forms, and acts on the block holding the quantity it constrains: a derivative's zero row is $\Brow(x^\ast)\G^{j}$ in the standard form and $\Brow(x^\ast)$ on the $j$-th derivative block of the descriptor's chain. $x^\ast$, the point or slice of the data; $x_s$, the datum point, where $f(x_s)\ne0$; $\tau(x)$, the Chebyshev polynomials evaluated at $x$; $\Slice_a$, slice along axis $a$; $\hat g$, coefficients of $g$; $a$, the anchor component of a vector datum $y$ and $y_\perp$ the rest; $\M_r$, multiplication by the source $r$; $\G$, differentiation matrix. Definitions: \Cref{sec:standard-constraints,sec:systems-data}.}
    \label{tab:constraint-kinds}
    \small
    \begin{tabular}{@{}l p{0.27\linewidth} p{0.34\linewidth} p{0.17\linewidth}@{}}
        \toprule
        Row & Data it carries & Operator & Problems \\
        \midrule
        \multicolumn{4}{@{}l}{\textit{Vanishes on the solution and its multiples: fixes the shape, leaves the scale free}} \\[2pt]
        Zero & $f^{(j)}(x^\ast)=0$ at a point: Dirichlet ($j=0$), Neumann ($j=1$), or the same on a boundary of a \gls{pde} & $\Brow(x^\ast)\G^{j}$, with $\Brow(x)=\sqrt{N}\ket{0}\bra{\tau(x)}$ & R2--R6 \\
        Datum slice & an initial or boundary \emph{function} $g$ on the slice $x_a=x^\ast$ of a \gls{pde} & $\Slice_a(x^\ast)-\hat g\bra{\tau(x_s)}/f(x_s)$ & R5 \\
        Ratio & a vector datum $y$ at $x^\ast$: the initial values of a coupled system or of a Carleman lift & $\bigl(\Frat/y_a\bigr)\otimes\Brow(x^\ast)$, with $\Frat=y_a(I-\ket{a}\bra{a})-y_\perp\bra{a}$ fixing the direction of $y$ & R7--R10 \\
        \addlinespace
        \multicolumn{4}{@{}l}{\textit{Maps the solution to a constant, whatever its scale: ties data to the scale}} \\[2pt]
        Collapse & the regular datum $f(x_s)$, which fixes the scale after preparation; through it a source $r$, a non-zero point datum $y^\ast$, and the linear terms of the doubled space & $\Dzero(x_s)=\Brow(x_s)/f(x_s)$, entering as $-\M_r\Dzero(x_s)$ for a source and as $\Brow(x^\ast)\G^{j}-y^\ast\Dzero(x_s)$ for a datum & every problem (the scale); R4 (sources); R6 \\
        \bottomrule
    \end{tabular}
\end{table}

Every run is a problem of \Cref{tab:ladder} at a resolution, with the options of \Cref{tab:run-inputs}. The defaults are the configuration every comparison of \Cref{sec:results} uses; the alternatives are controls, labelled ablations, or the printed choices kept to test the published method.

\begin{table}[htbp]
    \centering
    \caption[The model's input options and the values used]{\textbf{The model's input options and the values used.} Every run is a problem of \Cref{tab:ladder} at a resolution with these options, at the default unless a result states otherwise. $n$, qubits per Chebyshev axis ($N=2^n$); $n_t$, qubits on a time axis; $M$, Galerkin modes; $\zeta$, block weight; $K$, Carleman order; $\varepsilon$, target suppression; $r$, ratio of the geometric product state.}
    \label{tab:run-inputs}
    \small
    \begin{tabular}{@{}p{0.2\linewidth} p{0.38\linewidth} p{0.34\linewidth}@{}}
        \toprule
        Option & Values & Default and use \\
        \midrule
        Resolution & $n$ qubits per Chebyshev axis, $N=2^n$; $n_t$ on a time axis; modes $M$ (R9, R10) & swept per problem (\Cref{tab:catalogue-linear,tab:catalogue-nonlinear}) \\
        Form & standard; descriptor & both, compared \\
        Descriptor layout & chain (order $\ge2$); mass form (first order) & by the equation's order; the chain on a first-order system as an ablation \\
        Block rescaling & $\zeta$ from \Cref{eq:zeta-rule}; none & $\zeta$; unscaled as an ablation \\
        Encoding of $\G$ & structured, exact; Wu et al.'s & structured; Wu et al.'s in resource estimates only \\
        Nonlinearity & doubled space; Carleman order $K$, symmetric or tensor layout & as the problem requires; the doubled space in the standard form only \\
        Filter & minimax edge filter; imaginary time & minimax; imaginary time in the filter study and the reproduction \\
        Target $\varepsilon$ & suppression of the excited spectrum relative to the ground state & $10^{-6}$; also $10^{-4}$ and $10^{-8}$ in the filter study \\
        Initial state & geometric product state, $r=\pm\tfrac12$; uniform; random & geometric; uniform and random as controls \\
        Tier & circuit simulation (T1); exact emulation (T2); resource estimate (T3) & every tier the run's size allows (\Cref{sec:pipeline-verification}) \\
        \bottomrule
    \end{tabular}
\end{table}

\subsection{Results by problem}
\label{app:catalogue-results}

The linear benchmark problems and the coupled oscillator R7 in both forms, with the minimax edge filter at $\varepsilon=10^{-6}$ and the geometric initial state, under the conventions of \Cref{sec:pipeline-design}. \Cref{tab:catalogue-results} gives the largest resolution at which each form's reference is trusted and its preparation was emulated, with the prepared state's infidelity against the reference and the filter's degree. Every one of these preparations meets $\varepsilon$, its a priori bound and its field budget (\Cref{sec:results-faithful}). The standard form's emulations stop at smaller $n$ than the descriptor's because its degree grows as $N^{2k}$; a dash marks a problem with no trusted reference at the resolutions emulated, R2b and R2c in the standard form and the wave equation R5c in both, whose references are trusted only at larger $n$. Below $n=5$ per axis the wave equation's near-kernel is several-dimensional, $\res{build/R5c/descriptor/n4x4}{kernel_dimension}$ at $n=4$, consistent with a solution, $\cos4\pi t\,\sin2\pi x$, not yet resolved in $t$, and no preparation is claimed there. \Cref{tab:catalogue-t3} gives the resource estimates at the largest resolution estimated for every problem, as built and at the ideal-encoding bound $\alpha_R=\|R\|_2$, and \Cref{tab:catalogue-growth} the growth of these degrees with $N$, of which \Cref{tab:growth} carries a representative of each class.

\begin{table}[htbp]
    \centering
    \caption[Infidelity and degree of every trusted preparation]{\textbf{Every trusted preparation meets its target, in both forms.} Per problem and form: the largest resolution $n$ (per axis on R5) with a trusted reference and an exact emulation, the prepared state's infidelity against the reference, and the minimax filter's degree. --, no trusted reference at the resolutions emulated. Exact emulation (T2).}
    \label{tab:catalogue-results}
    \small
    \begin{tabular}{@{}l r r r r r r@{}}
        \toprule
        & \multicolumn{3}{c}{Standard} & \multicolumn{3}{c}{Descriptor} \\
        \cmidrule(lr){2-4}\cmidrule(lr){5-7}
        Problem & $n$ & Infidelity & Degree & $n$ & Infidelity & Degree \\
        \midrule
        R1 & $5$ & $\res[2,sci]{stateprep/R1/standard/n5/T2/filter=minimax}{infidelity}$ & $\res{stateprep/R1/standard/n5/T2/filter=minimax}{degree}$ & $8$ & $\res[2,sci]{stateprep/R1/descriptor/n8/T2/filter=minimax}{infidelity}$ & $\res{stateprep/R1/descriptor/n8/T2/filter=minimax}{degree}$ \\
        R2a & $3$ & $\res[2,sci]{stateprep/R2a/standard/n3/T2/filter=minimax}{infidelity}$ & $\res{stateprep/R2a/standard/n3/T2/filter=minimax}{degree}$ & $8$ & $\res[2,sci]{stateprep/R2a/descriptor/n8/T2/filter=minimax}{infidelity}$ & $\res{stateprep/R2a/descriptor/n8/T2/filter=minimax}{degree}$ \\
        R2b & \multicolumn{3}{c}{--} & $8$ & $\res[2,sci]{stateprep/R2b/descriptor/n8/T2/filter=minimax}{infidelity}$ & $\res{stateprep/R2b/descriptor/n8/T2/filter=minimax}{degree}$ \\
        R2c & \multicolumn{3}{c}{--} & $8$ & $\res[2,sci]{stateprep/R2c/descriptor/n8/T2/filter=minimax}{infidelity}$ & $\res{stateprep/R2c/descriptor/n8/T2/filter=minimax}{degree}$ \\
        R3a & $4$ & $\res[2,sci]{stateprep/R3a/standard/n4/T2/filter=minimax}{infidelity}$ & $\res{stateprep/R3a/standard/n4/T2/filter=minimax}{degree}$ & $7$ & $\res[2,sci]{stateprep/R3a/descriptor/n7/T2/filter=minimax}{infidelity}$ & $\res{stateprep/R3a/descriptor/n7/T2/filter=minimax}{degree}$ \\
        R3b & $4$ & $\res[2,sci]{stateprep/R3b/standard/n4/T2/filter=minimax}{infidelity}$ & $\res{stateprep/R3b/standard/n4/T2/filter=minimax}{degree}$ & $7$ & $\res[2,sci]{stateprep/R3b/descriptor/n7/T2/filter=minimax}{infidelity}$ & $\res{stateprep/R3b/descriptor/n7/T2/filter=minimax}{degree}$ \\
        R3c & $4$ & $\res[2,sci]{stateprep/R3c/standard/n4/T2/filter=minimax}{infidelity}$ & $\res{stateprep/R3c/standard/n4/T2/filter=minimax}{degree}$ & $7$ & $\res[2,sci]{stateprep/R3c/descriptor/n7/T2/filter=minimax}{infidelity}$ & $\res{stateprep/R3c/descriptor/n7/T2/filter=minimax}{degree}$ \\
        R3d & $4$ & $\res[2,sci]{stateprep/R3d/standard/n4/T2/filter=minimax}{infidelity}$ & $\res{stateprep/R3d/standard/n4/T2/filter=minimax}{degree}$ & $6$ & $\res[2,sci]{stateprep/R3d/descriptor/n6/T2/filter=minimax}{infidelity}$ & $\res{stateprep/R3d/descriptor/n6/T2/filter=minimax}{degree}$ \\
        R4a & $3$ & $\res[2,sci]{stateprep/R4a/standard/n3/T2/filter=minimax}{infidelity}$ & $\res{stateprep/R4a/standard/n3/T2/filter=minimax}{degree}$ & $7$ & $\res[2,sci]{stateprep/R4a/descriptor/n7/T2/filter=minimax}{infidelity}$ & $\res{stateprep/R4a/descriptor/n7/T2/filter=minimax}{degree}$ \\
        R4b & $4$ & $\res[2,sci]{stateprep/R4b/standard/n4/T2/filter=minimax}{infidelity}$ & $\res{stateprep/R4b/standard/n4/T2/filter=minimax}{degree}$ & $8$ & $\res[2,sci]{stateprep/R4b/descriptor/n8/T2/filter=minimax}{infidelity}$ & $\res{stateprep/R4b/descriptor/n8/T2/filter=minimax}{degree}$ \\
        R4c & $4$ & $\res[2,sci]{stateprep/R4c/standard/n4/T2/filter=minimax}{infidelity}$ & $\res{stateprep/R4c/standard/n4/T2/filter=minimax}{degree}$ & $8$ & $\res[2,sci]{stateprep/R4c/descriptor/n8/T2/filter=minimax}{infidelity}$ & $\res{stateprep/R4c/descriptor/n8/T2/filter=minimax}{degree}$ \\
        R4d & $4$ & $\res[2,sci]{stateprep/R4d/standard/n4/T2/filter=minimax}{infidelity}$ & $\res{stateprep/R4d/standard/n4/T2/filter=minimax}{degree}$ & $8$ & $\res[2,sci]{stateprep/R4d/descriptor/n8/T2/filter=minimax}{infidelity}$ & $\res{stateprep/R4d/descriptor/n8/T2/filter=minimax}{degree}$ \\
        R5a & $3$ & $\res[2,sci]{stateprep/R5a/standard/n3x3/T2/filter=minimax}{infidelity}$ & $\res{stateprep/R5a/standard/n3x3/T2/filter=minimax}{degree}$ & $4$ & $\res[2,sci]{stateprep/R5a/descriptor/n4x4/T2/filter=minimax}{infidelity}$ & $\res{stateprep/R5a/descriptor/n4x4/T2/filter=minimax}{degree}$ \\
        R5b & $4$ & $\res[2,sci]{stateprep/R5b/standard/n4x4/T2/filter=minimax}{infidelity}$ & $\res{stateprep/R5b/standard/n4x4/T2/filter=minimax}{degree}$ & $5$ & $\res[2,sci]{stateprep/R5b/descriptor/n5x5/T2/filter=minimax}{infidelity}$ & $\res{stateprep/R5b/descriptor/n5x5/T2/filter=minimax}{degree}$ \\
        R5c & \multicolumn{3}{c}{--} & \multicolumn{3}{c}{--} \\
        R7 & $7$ & $\res[2,sci]{stateprep/R7/standard/n7/T2/filter=minimax}{infidelity}$ & $\res{stateprep/R7/standard/n7/T2/filter=minimax}{degree}$ & $9$ & $\res[2,sci]{stateprep/R7/descriptor/n9/T2/filter=minimax}{infidelity}$ & $\res{stateprep/R7/descriptor/n9/T2/filter=minimax}{degree}$ \\
        \bottomrule
    \end{tabular}
\end{table}

\begin{table}[htbp]
    \centering
    \caption[Filter degree at the largest resolution estimated, every problem]{\textbf{As built the descriptor's degree is the lower on every problem; at the ideal-encoding bound the standard form's is the lower on R3.} Minimax edge filter degree at $\varepsilon=10^{-6}$, $n=8$ on one axis and $n=5$ per axis on R5, as built and at the ideal-encoding bound $\alpha_R=\|R\|_2$ ($\alpha_R$, the residual's subnormalisation; $\|R\|_2$, its norm). At the ideal bound the standard form's degree is the lower on R3a--R3d and the two are close on R5b. Resource estimates (T3).}
    \label{tab:catalogue-t3}
    \small
    \begin{tabular}{@{}l r r r r r@{}}
        \toprule
        & & \multicolumn{2}{c}{As built} & \multicolumn{2}{c}{$\alpha_R=\|R\|_2$} \\
        \cmidrule(lr){3-4}\cmidrule(lr){5-6}
        Problem & $n$ & Standard & Descriptor & Standard & Descriptor \\
        \midrule
        R1 & $8$ & $\res[2,sci]{stateprep/R1/standard/n8/T3/filter=minimax}{degree}$ & $\res[2,sci]{stateprep/R1/descriptor/n8/T3/filter=minimax}{degree}$ & $\res[2,sci]{stateprep/R1/standard/n8/T3/filter=minimax}{ideal_degree}$ & $\res[2,sci]{stateprep/R1/descriptor/n8/T3/filter=minimax}{ideal_degree}$ \\
        R2a & $8$ & $\res[2,sci]{stateprep/R2a/standard/n8/T3/filter=minimax}{degree}$ & $\res[2,sci]{stateprep/R2a/descriptor/n8/T3/filter=minimax}{degree}$ & $\res[2,sci]{stateprep/R2a/standard/n8/T3/filter=minimax}{ideal_degree}$ & $\res[2,sci]{stateprep/R2a/descriptor/n8/T3/filter=minimax}{ideal_degree}$ \\
        R2b & $8$ & $\res[2,sci]{stateprep/R2b/standard/n8/T3/filter=minimax}{degree}$ & $\res[2,sci]{stateprep/R2b/descriptor/n8/T3/filter=minimax}{degree}$ & $\res[2,sci]{stateprep/R2b/standard/n8/T3/filter=minimax}{ideal_degree}$ & $\res[2,sci]{stateprep/R2b/descriptor/n8/T3/filter=minimax}{ideal_degree}$ \\
        R2c & $8$ & $\res[2,sci]{stateprep/R2c/standard/n8/T3/filter=minimax}{degree}$ & $\res[2,sci]{stateprep/R2c/descriptor/n8/T3/filter=minimax}{degree}$ & $\res[2,sci]{stateprep/R2c/standard/n8/T3/filter=minimax}{ideal_degree}$ & $\res[2,sci]{stateprep/R2c/descriptor/n8/T3/filter=minimax}{ideal_degree}$ \\
        R3a & $8$ & $\res[2,sci]{stateprep/R3a/standard/n8/T3/filter=minimax}{degree}$ & $\res[2,sci]{stateprep/R3a/descriptor/n8/T3/filter=minimax}{degree}$ & $\res[2,sci]{stateprep/R3a/standard/n8/T3/filter=minimax}{ideal_degree}$ & $\res[2,sci]{stateprep/R3a/descriptor/n8/T3/filter=minimax}{ideal_degree}$ \\
        R3b & $8$ & $\res[2,sci]{stateprep/R3b/standard/n8/T3/filter=minimax}{degree}$ & $\res[2,sci]{stateprep/R3b/descriptor/n8/T3/filter=minimax}{degree}$ & $\res[2,sci]{stateprep/R3b/standard/n8/T3/filter=minimax}{ideal_degree}$ & $\res[2,sci]{stateprep/R3b/descriptor/n8/T3/filter=minimax}{ideal_degree}$ \\
        R3c & $8$ & $\res[2,sci]{stateprep/R3c/standard/n8/T3/filter=minimax}{degree}$ & $\res[2,sci]{stateprep/R3c/descriptor/n8/T3/filter=minimax}{degree}$ & $\res[2,sci]{stateprep/R3c/standard/n8/T3/filter=minimax}{ideal_degree}$ & $\res[2,sci]{stateprep/R3c/descriptor/n8/T3/filter=minimax}{ideal_degree}$ \\
        R3d & $8$ & $\res[2,sci]{stateprep/R3d/standard/n8/T3/filter=minimax}{degree}$ & $\res[2,sci]{stateprep/R3d/descriptor/n8/T3/filter=minimax}{degree}$ & $\res[2,sci]{stateprep/R3d/standard/n8/T3/filter=minimax}{ideal_degree}$ & $\res[2,sci]{stateprep/R3d/descriptor/n8/T3/filter=minimax}{ideal_degree}$ \\
        R4a & $8$ & $\res[2,sci]{stateprep/R4a/standard/n8/T3/filter=minimax}{degree}$ & $\res[2,sci]{stateprep/R4a/descriptor/n8/T3/filter=minimax}{degree}$ & $\res[2,sci]{stateprep/R4a/standard/n8/T3/filter=minimax}{ideal_degree}$ & $\res[2,sci]{stateprep/R4a/descriptor/n8/T3/filter=minimax}{ideal_degree}$ \\
        R4b & $8$ & $\res[2,sci]{stateprep/R4b/standard/n8/T3/filter=minimax}{degree}$ & $\res[2,sci]{stateprep/R4b/descriptor/n8/T3/filter=minimax}{degree}$ & $\res[2,sci]{stateprep/R4b/standard/n8/T3/filter=minimax}{ideal_degree}$ & $\res[2,sci]{stateprep/R4b/descriptor/n8/T3/filter=minimax}{ideal_degree}$ \\
        R4c & $8$ & $\res[2,sci]{stateprep/R4c/standard/n8/T3/filter=minimax}{degree}$ & $\res[2,sci]{stateprep/R4c/descriptor/n8/T3/filter=minimax}{degree}$ & $\res[2,sci]{stateprep/R4c/standard/n8/T3/filter=minimax}{ideal_degree}$ & $\res[2,sci]{stateprep/R4c/descriptor/n8/T3/filter=minimax}{ideal_degree}$ \\
        R4d & $8$ & $\res[2,sci]{stateprep/R4d/standard/n8/T3/filter=minimax}{degree}$ & $\res[2,sci]{stateprep/R4d/descriptor/n8/T3/filter=minimax}{degree}$ & $\res[2,sci]{stateprep/R4d/standard/n8/T3/filter=minimax}{ideal_degree}$ & $\res[2,sci]{stateprep/R4d/descriptor/n8/T3/filter=minimax}{ideal_degree}$ \\
        R5a & $5$ & $\res[2,sci]{stateprep/R5a/standard/n5x5/T3/filter=minimax}{degree}$ & $\res[2,sci]{stateprep/R5a/descriptor/n5x5/T3/filter=minimax}{degree}$ & $\res[2,sci]{stateprep/R5a/standard/n5x5/T3/filter=minimax}{ideal_degree}$ & $\res[2,sci]{stateprep/R5a/descriptor/n5x5/T3/filter=minimax}{ideal_degree}$ \\
        R5b & $5$ & $\res[2,sci]{stateprep/R5b/standard/n5x5/T3/filter=minimax}{degree}$ & $\res[2,sci]{stateprep/R5b/descriptor/n5x5/T3/filter=minimax}{degree}$ & $\res[2,sci]{stateprep/R5b/standard/n5x5/T3/filter=minimax}{ideal_degree}$ & $\res[2,sci]{stateprep/R5b/descriptor/n5x5/T3/filter=minimax}{ideal_degree}$ \\
        R5c & $5$ & $\res[2,sci]{stateprep/R5c/standard/n5x5/T3/filter=minimax}{degree}$ & $\res[2,sci]{stateprep/R5c/descriptor/n5x5/T3/filter=minimax}{degree}$ & $\res[2,sci]{stateprep/R5c/standard/n5x5/T3/filter=minimax}{ideal_degree}$ & $\res[2,sci]{stateprep/R5c/descriptor/n5x5/T3/filter=minimax}{ideal_degree}$ \\
        R7 & $8$ & $\res[2,sci]{stateprep/R7/standard/n8/T3/filter=minimax}{degree}$ & $\res[2,sci]{stateprep/R7/descriptor/n8/T3/filter=minimax}{degree}$ & $\res[2,sci]{stateprep/R7/standard/n8/T3/filter=minimax}{ideal_degree}$ & $\res[2,sci]{stateprep/R7/descriptor/n8/T3/filter=minimax}{ideal_degree}$ \\
        \bottomrule
    \end{tabular}
\end{table}

\begin{table}[htbp]
    \centering
    \caption[Growth of the filter's degree with $N$, every problem]{\textbf{The degree's growth with $N$ on every linear problem and on R7.} Exponent of $N=2^n$ in the minimax filter's degree between the two largest resolutions estimated ($n=7\to8$ on one axis, $n=4\to5$ per axis on two), as built and at the ideal-encoding bound $\alpha_R=\|R\|_2$, grouped by class: $k$, equation order; regular, singular: the leading coefficient without and with a zero on $[-1,1]$ (\Cref{def:regular}). $^\dagger$Reference not trusted at $n=4$ per axis; slope indicative only. Two-point slopes of a degree with logarithmic factors (\Cref{sec:results-scaling}). Resource estimates (T3).}
    \label{tab:catalogue-growth}
    \small
    \begin{tabular}{@{}l l r r r r@{}}
        \toprule
        & & \multicolumn{2}{c}{Standard} & \multicolumn{2}{c}{Descriptor} \\
        \cmidrule(lr){3-4}\cmidrule(lr){5-6}
        Class & Problem & As built & Ideal & As built & Ideal \\
        \midrule
        First order, $k=1$ & R1 & $\res[3]{stateprep/R1/standard/growth/T3/filter=minimax}{degree}$ & $\res[3]{stateprep/R1/standard/growth/T3/filter=minimax}{ideal_degree}$ & $\res[3]{stateprep/R1/descriptor/growth/T3/filter=minimax}{degree}$ & $\res[3]{stateprep/R1/descriptor/growth/T3/filter=minimax}{ideal_degree}$ \\
         & R7 & $\res[3]{stateprep/R7/standard/growth/T3/filter=minimax}{degree}$ & $\res[3]{stateprep/R7/standard/growth/T3/filter=minimax}{ideal_degree}$ & $\res[3]{stateprep/R7/descriptor/growth/T3/filter=minimax}{degree}$ & $\res[3]{stateprep/R7/descriptor/growth/T3/filter=minimax}{ideal_degree}$ \\
        \addlinespace
        Second order, regular & R2a & $\res[3]{stateprep/R2a/standard/growth/T3/filter=minimax}{degree}$ & $\res[3]{stateprep/R2a/standard/growth/T3/filter=minimax}{ideal_degree}$ & $\res[3]{stateprep/R2a/descriptor/growth/T3/filter=minimax}{degree}$ & $\res[3]{stateprep/R2a/descriptor/growth/T3/filter=minimax}{ideal_degree}$ \\
         & R2b & $\res[3]{stateprep/R2b/standard/growth/T3/filter=minimax}{degree}$ & $\res[3]{stateprep/R2b/standard/growth/T3/filter=minimax}{ideal_degree}$ & $\res[3]{stateprep/R2b/descriptor/growth/T3/filter=minimax}{degree}$ & $\res[3]{stateprep/R2b/descriptor/growth/T3/filter=minimax}{ideal_degree}$ \\
         & R2c & $\res[3]{stateprep/R2c/standard/growth/T3/filter=minimax}{degree}$ & $\res[3]{stateprep/R2c/standard/growth/T3/filter=minimax}{ideal_degree}$ & $\res[3]{stateprep/R2c/descriptor/growth/T3/filter=minimax}{degree}$ & $\res[3]{stateprep/R2c/descriptor/growth/T3/filter=minimax}{ideal_degree}$ \\
         & R4b & $\res[3]{stateprep/R4b/standard/growth/T3/filter=minimax}{degree}$ & $\res[3]{stateprep/R4b/standard/growth/T3/filter=minimax}{ideal_degree}$ & $\res[3]{stateprep/R4b/descriptor/growth/T3/filter=minimax}{degree}$ & $\res[3]{stateprep/R4b/descriptor/growth/T3/filter=minimax}{ideal_degree}$ \\
         & R4c & $\res[3]{stateprep/R4c/standard/growth/T3/filter=minimax}{degree}$ & $\res[3]{stateprep/R4c/standard/growth/T3/filter=minimax}{ideal_degree}$ & $\res[3]{stateprep/R4c/descriptor/growth/T3/filter=minimax}{degree}$ & $\res[3]{stateprep/R4c/descriptor/growth/T3/filter=minimax}{ideal_degree}$ \\
         & R4d & $\res[3]{stateprep/R4d/standard/growth/T3/filter=minimax}{degree}$ & $\res[3]{stateprep/R4d/standard/growth/T3/filter=minimax}{ideal_degree}$ & $\res[3]{stateprep/R4d/descriptor/growth/T3/filter=minimax}{degree}$ & $\res[3]{stateprep/R4d/descriptor/growth/T3/filter=minimax}{ideal_degree}$ \\
        \addlinespace
        Second order, singular & R3a & $\res[3]{stateprep/R3a/standard/growth/T3/filter=minimax}{degree}$ & $\res[3]{stateprep/R3a/standard/growth/T3/filter=minimax}{ideal_degree}$ & $\res[3]{stateprep/R3a/descriptor/growth/T3/filter=minimax}{degree}$ & $\res[3]{stateprep/R3a/descriptor/growth/T3/filter=minimax}{ideal_degree}$ \\
         & R3b & $\res[3]{stateprep/R3b/standard/growth/T3/filter=minimax}{degree}$ & $\res[3]{stateprep/R3b/standard/growth/T3/filter=minimax}{ideal_degree}$ & $\res[3]{stateprep/R3b/descriptor/growth/T3/filter=minimax}{degree}$ & $\res[3]{stateprep/R3b/descriptor/growth/T3/filter=minimax}{ideal_degree}$ \\
         & R3c & $\res[3]{stateprep/R3c/standard/growth/T3/filter=minimax}{degree}$ & $\res[3]{stateprep/R3c/standard/growth/T3/filter=minimax}{ideal_degree}$ & $\res[3]{stateprep/R3c/descriptor/growth/T3/filter=minimax}{degree}$ & $\res[3]{stateprep/R3c/descriptor/growth/T3/filter=minimax}{ideal_degree}$ \\
         & R3d & $\res[3]{stateprep/R3d/standard/growth/T3/filter=minimax}{degree}$ & $\res[3]{stateprep/R3d/standard/growth/T3/filter=minimax}{ideal_degree}$ & $\res[3]{stateprep/R3d/descriptor/growth/T3/filter=minimax}{degree}$ & $\res[3]{stateprep/R3d/descriptor/growth/T3/filter=minimax}{ideal_degree}$ \\
         & R4a & $\res[3]{stateprep/R4a/standard/growth/T3/filter=minimax}{degree}$ & $\res[3]{stateprep/R4a/standard/growth/T3/filter=minimax}{ideal_degree}$ & $\res[3]{stateprep/R4a/descriptor/growth/T3/filter=minimax}{degree}$ & $\res[3]{stateprep/R4a/descriptor/growth/T3/filter=minimax}{ideal_degree}$ \\
        \addlinespace
        Two axes, per axis & R5a & $\res[3]{stateprep/R5a/standard/growth/T3/filter=minimax}{degree}$ & $\res[3]{stateprep/R5a/standard/growth/T3/filter=minimax}{ideal_degree}$ & $\res[3]{stateprep/R5a/descriptor/growth/T3/filter=minimax}{degree}$ & $\res[3]{stateprep/R5a/descriptor/growth/T3/filter=minimax}{ideal_degree}$ \\
         & R5b & $\res[3]{stateprep/R5b/standard/growth/T3/filter=minimax}{degree}$ & $\res[3]{stateprep/R5b/standard/growth/T3/filter=minimax}{ideal_degree}$ & $\res[3]{stateprep/R5b/descriptor/growth/T3/filter=minimax}{degree}$ & $\res[3]{stateprep/R5b/descriptor/growth/T3/filter=minimax}{ideal_degree}$ \\
         & R5c$^\dagger$ & $\res[3]{stateprep/R5c/standard/growth/T3/filter=minimax}{degree}$ & $\res[3]{stateprep/R5c/standard/growth/T3/filter=minimax}{ideal_degree}$ & $\res[3]{stateprep/R5c/descriptor/growth/T3/filter=minimax}{degree}$ & $\res[3]{stateprep/R5c/descriptor/growth/T3/filter=minimax}{ideal_degree}$ \\
        \bottomrule
    \end{tabular}
\end{table}

\subsection{Filter studies}
\label{app:filter-studies}

\paragraph{Comparing the two filters at like error.}
\label{app:filter-study}
A degree means something only at a stated error, so the filter study of \Cref{sec:results-faithful} holds both filters to the same one, at three strengths. \emph{Equal suppression} $\eta$ bounds the excited interval relative to the ground state, $\sup_{x\in[x_1,1]}|P(x)|/|P(x_0)|\le\eta$; it bounds every initial state's infidelity by $\eta^2(1-\gamma^2)/\gamma^2$, depends only on $x_0$ and $x_1$, and is the strength at which the minimax filter is optimal (\Cref{prop:minimax-optimal}). \emph{Equal certified infidelity} $\varepsilon$ is the same bound at the problem's own overlap, $\eta=\gamma\sqrt{\varepsilon/(1-\gamma^2)}$. \emph{Equal measured infidelity} $\varepsilon$ is reached on the problem's own spectrum and initial state, and a degree counts only if every larger degree up to the certified one also reaches $\varepsilon$, so that neither filter's oscillation lands on $\varepsilon$ by coincidence. Imaginary time's $\beta$ is free: each of its degrees is the least over $\beta$, found by a grid search refined about its best point, and so an upper bound on the true least to the refinement's spacing. Both filters are anchored at the spectrum's edge with unit weight (\Cref{prop:filter-normalisation}), so \gls{gqsp} realises either at the same rescaling, and a filter whose success probability falls below $10^{-12}$ is not counted.

On synthetic spectra with the ground state at the edge, gaps from $10^{-8}$ to $10^{-1}$ on the scale of $H'$ and suppressions from $10^{-2}$ to $10^{-10}$, imaginary time needs $\res[3]{filters/synthetic/edge}{low}$--$\res[3]{filters/synthetic/edge}{high}$ times the minimax filter's degree at equal suppression, varying with the suppression and hardly with the gap; above the edge the ratio grows, since imaginary time's value on the ground state then falls with its time. \Cref{tab:filter-study} gives the study on the problems whose references both forms trust. At equal certified infidelity both filters' degrees come from the same worst-case bound, so the ratio is the equal-suppression ratio at the problem's own $\gamma$; every trusted reference has its ground state at the edge, and the ratio lies within the synthetic range in both forms.
It does not change with the resolution: on R1 in the descriptor it is $\res{filters/R1/descriptor/n2}{e6_qite_worst}$ / $\res{filters/R1/descriptor/n2}{e6_minimax_certified}$ at $n=2$ and $\res{filters/R1/descriptor/n8}{e6_qite_worst}$ / $\res{filters/R1/descriptor/n8}{e6_minimax_certified}$ at $n=8$, so the filter changes the degree's constant and not its growth with $N$. At equal measured infidelity the ratio is smaller, because imaginary time suppresses each excited state more the higher it lies, while the bound holds every one at the equiripple level. At $\varepsilon=10^{-4}$ imaginary time needs fewer degrees on R1 in the standard form and on R3a and R5a in the descriptor, and the same on R1 in the descriptor; at $10^{-6}$ the minimax filter needs fewer on every case.

\begin{table}[htbp]
    \centering
    \caption[Filter degree, imaginary time against minimax]{\textbf{The minimax edge filter needs a lower degree than imaginary time at equal error, except at loose $\varepsilon$ on some problems.} Each cell: imaginary-time (QITE) degree / minimax degree, both forms, at one resolution per problem ($n$ per axis on R5a) whose reference both forms trust. Three equal-error criteria: certified infidelity at $\varepsilon=10^{-6}$ (both degrees from the worst-case bound at the problem's own $\gamma$), and measured infidelity at $\varepsilon=10^{-6}$ and $10^{-4}$ on the problem's own spectrum from the geometric initial state. $\gamma^2$, that state's overlap with the ground state. Exact emulation (T2).}
    \label{tab:filter-study}
    \footnotesize
    \setlength{\tabcolsep}{4pt}
    \begin{tabular}{@{}l l r r r r@{}}
        \toprule
        Problem & Form & $\gamma^2$ & Certified, $10^{-6}$ & Measured, $10^{-6}$ & Measured, $10^{-4}$ \\
        \midrule
        R1, $n=4$ & standard & $\res[3]{filters/R1/standard/n4}{gamma2}$ & $\res{filters/R1/standard/n4}{e6_qite_worst}$ / $\res{filters/R1/standard/n4}{e6_minimax_certified}$ & $\res{filters/R1/standard/n4}{e6_qite_degree}$ / $\res{filters/R1/standard/n4}{e6_minimax_degree}$ & $\res{filters/R1/standard/n4}{e4_qite_degree}$ / $\res{filters/R1/standard/n4}{e4_minimax_degree}$ \\
         & descriptor & $\res[3]{filters/R1/descriptor/n4}{gamma2}$ & $\res{filters/R1/descriptor/n4}{e6_qite_worst}$ / $\res{filters/R1/descriptor/n4}{e6_minimax_certified}$ & $\res{filters/R1/descriptor/n4}{e6_qite_degree}$ / $\res{filters/R1/descriptor/n4}{e6_minimax_degree}$ & $\res{filters/R1/descriptor/n4}{e4_qite_degree}$ / $\res{filters/R1/descriptor/n4}{e4_minimax_degree}$ \\
        R2a, $n=3$ & standard & $\res[3]{filters/R2a/standard/n3}{gamma2}$ & $\res{filters/R2a/standard/n3}{e6_qite_worst}$ / $\res{filters/R2a/standard/n3}{e6_minimax_certified}$ & $\res{filters/R2a/standard/n3}{e6_qite_degree}$ / $\res{filters/R2a/standard/n3}{e6_minimax_degree}$ & $\res{filters/R2a/standard/n3}{e4_qite_degree}$ / $\res{filters/R2a/standard/n3}{e4_minimax_degree}$ \\
         & descriptor & $\res[3]{filters/R2a/descriptor/n3}{gamma2}$ & $\res{filters/R2a/descriptor/n3}{e6_qite_worst}$ / $\res{filters/R2a/descriptor/n3}{e6_minimax_certified}$ & $\res{filters/R2a/descriptor/n3}{e6_qite_degree}$ / $\res{filters/R2a/descriptor/n3}{e6_minimax_degree}$ & $\res{filters/R2a/descriptor/n3}{e4_qite_degree}$ / $\res{filters/R2a/descriptor/n3}{e4_minimax_degree}$ \\
        R3a, $n=3$ & standard & $\res[3]{filters/R3a/standard/n3}{gamma2}$ & $\res{filters/R3a/standard/n3}{e6_qite_worst}$ / $\res{filters/R3a/standard/n3}{e6_minimax_certified}$ & $\res{filters/R3a/standard/n3}{e6_qite_degree}$ / $\res{filters/R3a/standard/n3}{e6_minimax_degree}$ & $\res{filters/R3a/standard/n3}{e4_qite_degree}$ / $\res{filters/R3a/standard/n3}{e4_minimax_degree}$ \\
         & descriptor & $\res[2,sci]{filters/R3a/descriptor/n3}{gamma2}$ & $\res{filters/R3a/descriptor/n3}{e6_qite_worst}$ / $\res{filters/R3a/descriptor/n3}{e6_minimax_certified}$ & $\res{filters/R3a/descriptor/n3}{e6_qite_degree}$ / $\res{filters/R3a/descriptor/n3}{e6_minimax_degree}$ & $\res{filters/R3a/descriptor/n3}{e4_qite_degree}$ / $\res{filters/R3a/descriptor/n3}{e4_minimax_degree}$ \\
        R5a, $n=3$ & standard & $\res[3]{filters/R5a/standard/n3x3}{gamma2}$ & $\res{filters/R5a/standard/n3x3}{e6_qite_worst}$ / $\res{filters/R5a/standard/n3x3}{e6_minimax_certified}$ & $\res{filters/R5a/standard/n3x3}{e6_qite_degree}$ / $\res{filters/R5a/standard/n3x3}{e6_minimax_degree}$ & $\res{filters/R5a/standard/n3x3}{e4_qite_degree}$ / $\res{filters/R5a/standard/n3x3}{e4_minimax_degree}$ \\
         & descriptor & $\res[3]{filters/R5a/descriptor/n3x3}{gamma2}$ & $\res{filters/R5a/descriptor/n3x3}{e6_qite_worst}$ / $\res{filters/R5a/descriptor/n3x3}{e6_minimax_certified}$ & $\res{filters/R5a/descriptor/n3x3}{e6_qite_degree}$ / $\res{filters/R5a/descriptor/n3x3}{e6_minimax_degree}$ & $\res{filters/R5a/descriptor/n3x3}{e4_qite_degree}$ / $\res{filters/R5a/descriptor/n3x3}{e4_minimax_degree}$ \\
        \bottomrule
    \end{tabular}
\end{table}

\paragraph{An underestimated edge.}
\label{app:edge-study}
The degree is set a priori from the spectral edge (\Cref{sec:pipeline-filter}), and a study on R2a at $n=3$ in the standard form measures what an underestimate of it costs, with the imaginary-time filter, whose time the gap sets directly (\Cref{eq:qite-parameters}). At the true gap the degree is $\res{protocol/R2a/standard/n3}{gap_degree}$ and the infidelity $\res[2,sci]{protocol/R2a/standard/n3}{gap_infidelity}$. With the gap halved the degree is $\res{protocol/R2a/standard/n3}{gap_half_degree}$, and with it divided by ten $\res{protocol/R2a/standard/n3}{gap_tenth_degree}$: the factors $\sqrt2$ and $\sqrt{10}$ of an inverse square root, with every preparation still within $\varepsilon$. Wu et al.'s own time rule, run in the same study, gives degree $\res{protocol/R2a/standard/n3}{paper_degree}$ and infidelity $\res[2]{protocol/R2a/standard/n3}{paper_infidelity}$: its time is shorter than the gap rule's and its state does not reach $\varepsilon=10^{-6}$. The time rule is examined with the rest of the reproduction in \Cref{app:reproduction}. Exact emulation (T2).
